% Ampliación demostrativa DF003--DF007. No modifica la fuente integral.
% Propietario común de los archivos Lean citados en estos comentarios:
% output/PAQUETE_ARTICULO_I_INTEGRACION_ESTADO_CAMPOS_20260919/.
% Residencia causal de la primera sección: después de los lectores regionales
% coinductivos y antes de emplear sus registros como entrada de la selección K.
% Residencia causal de la segunda sección: después de la construcción marcada
% Paley--Witt--Golay--Leech y antes de la reconstrucción de la estructura VOA.
% Los dos cortes expositivos conservan el mismo origen APP--TRIT--TPK.

\section[Génération régionale et raffinement]{Compatibilité intégrale de la génération régionale, du raffinement cylindrique et de l'incidence}
\label{act21:sec:df-lectores-cilindros}

La prolongation régionale produit des suites ordonnées de blocs, ainsi que
leurs profondeurs de lecture. La projection archimédienne et l'incidence
ternaire se calculent sur ces mêmes suites. Le lien entre ces deux
réalisations peut s'exprimer entièrement par des opérations entières :
sélection par intervalles rationnels, division euclidienne, élévation de six
trits, lecture de retenue et transformation de Witt. Cette formulation permet
de suivre chaque coefficient depuis son bloc générateur jusqu'à son emploi ultérieur.

La construction utilise les trois lecteurs régionaux déjà obtenus de l'état APP--TRIT--TPK. Nous les désignerons par les indices \(c\in\{\mathrm P,\mathrm E,\mathrm F\}\), correspondant respectivement aux régions de clôture, de propagation et d'auto-échelle. Leur réalisation réelle s'écrira \(v_c\). Les identifications ultérieures \(v_{\mathrm P}=\pi_{\mathrm{HMT}}\), \(v_{\mathrm E}=e_{\mathrm{HMT}}\) et \(v_{\mathrm F}=\varphi_{\mathrm{HMT}}\) conservent cet ordre causal. L'opération qui émet un bloc inspecte les intervalles rationnels du lecteur régional ; la valeur réelle intervient dans la démonstration de correction de cette opération.

\subsection{Émission positionnelle par séparation rationnelle de cellules}
\label{act21:subsec:df-emision-racional}

% Fuente: lean/biblioteca/RegionalPublicationComposition.lean,
% produced_brackets, eventual_publication, joint_eventual_publication,
% search_terminates, stopped_cell_correct, publications_compatible.

Pour chaque région, on dispose de deux suites rationnelles
\(\ell_c(j),u_c(j)\), construites par son lecteur, qui satisfont
\[
 \ell_c(j)\leq v_c\leq u_c(j).
\]
La propriété de séparation des cellules démontrée pour ces lecteurs affirme
que, pour chaque entier \(s>0\), il existe \(J\) tel que, pour tout \(j\geq J\),
\begin{equation}
 \lfloor s\ell_c(j)\rfloor
 =\lceil su_c(j)\rceil-1
 =\lfloor sv_c\rfloor.
 \label{act21:eq:df-separacion-celdas}
\end{equation}
L'évitement des frontières rationnelles par les valeurs régionales est l'hypothèse des théorèmes précédents qui garantit cette séparation. Ici, on emploie sa conclusion opérationnelle, avec les lecteurs et les intervalles qui la produisent.

\begin{definition}[Indice de séparation et préfixe émis]
Pour \(s>0\), soit
\[
 j_c(s)=\min\{j\in\mathbb N:
 \lfloor s\ell_c(j)\rfloor=\lceil su_c(j)\rceil-1\}.
\]
Pour une base entière \(b\geq2\) et une profondeur \(n\geq0\), on définit
\begin{equation}
 P_c(b,n)=\lfloor b^n\ell_c(j_c(b^n))\rfloor.
 \label{act21:eq:df-publicacion-prefijo}
\end{equation}
Le préfixe contient la partie entière et les premières \(n\) positions en base \(b\) ; par conséquent, \(P_c(b,0)\) est la partie entière régionale.
\end{definition}

\begin{theorem}[Terminaison et correction de l'émission]
\label{act21:thm:df-emision-correcta}
L'indice \(j_c(s)\) existe pour tout \(s>0\). Les préfixes émis satisfont
\begin{equation}
 P_c(b,n)=\lfloor b^n v_c\rfloor,
 \qquad
 P_c(b,n)\leq b^n v_c<P_c(b,n)+1.
 \label{act21:eq:df-prefijo-correcto}
\end{equation}
Pour chaque échelle \(s\) il est possible en outre de choisir un unique indice de suffisance \(J(s)\) valide simultanément pour les trois régions.
\end{theorem}

\begin{proof}
L'égalité à partir d'un certain rang de \eqref{act21:eq:df-separacion-celdas} rend non vide
l'ensemble qui définit \(j_c(s)\). Son minimum existe par le bon ordre de
\(\mathbb N\). À tout indice où l'égalité se produit, la borne
inférieure rationnelle fournit
\(\lfloor s\ell_c(j)\rfloor\leq\lfloor sv_c\rfloor\).
La borne supérieure, jointe à l'évitement d'une frontière de cellule, donne
\(\lfloor sv_c\rfloor\leq\lceil su_c(j)\rceil-1\).
Les deux extrémités entières coïncidant, l'entier intermédiaire coïncide avec
elles. Cela démontre \eqref{act21:eq:df-prefijo-correcto}. Pour l'affirmation
conjointe, il suffit de prendre le maximum des trois indices de suffisance obtenus
pour la même échelle. L'indice conjoint contrôle simultanément les trois
émissions, tandis que chaque région conserve ses propres intervalles.
\end{proof}

\begin{proposition}[Compatibilité exacte des préfixes]
\label{act21:prop:df-prefijos-compatibles}
Pour \(n,m\geq0\),
\begin{equation}
 \left\lfloor\frac{P_c(b,n+m)}{b^m}\right\rfloor=P_c(b,n).
 \label{act21:eq:df-truncamiento}
\end{equation}
En particulier, le bloc suivant
\(d_c(b,n)=P_c(b,n+1)\bmod b\) satisfait
\begin{equation}
 P_c(b,n+1)=bP_c(b,n)+d_c(b,n),\qquad 0\leq d_c(b,n)<b.
 \label{act21:eq:df-paso-posicional}
\end{equation}
\end{proposition}

\begin{proof}
Écrivons \(b^n v_c=N+\theta\), avec
\(N=P_c(b,n)\) et \(0\leq\theta<1\). Alors
\(P_c(b,n+m)=b^mN+\lfloor b^m\theta\rfloor\), et le second terme
appartient à \(\{0,\ldots,b^m-1\}\). La division euclidienne par \(b^m\)
récupère \(N\). La spécialisation \(m=1\) produit
\eqref{act21:eq:df-paso-posicional} et sa borne résiduelle.
\end{proof}

\subsection{Blocs régionaux de six trits à profondeur arbitraire}
\label{act21:subsec:df-bloques-arbitrarios}

% Fuente: deltas/integracion/JointRegionalFrontier.lean,
% publish_base729_eq_base3, generated_block_from_longer_prefix,
% generated_block_eq_finite, generated_signature_eq_finite.

L'égalité entière \(729=3^6\) détermine le regroupement de six positions ternaires en un bloc. Pour \(t\geq0\), nous définissons
\begin{equation}
 u_c(t)=P_c(729,t+1)\bmod729,
 \qquad
 B_{c,j}(t)=
 \left\lfloor\frac{u_c(t)}{3^{5-j}}\right\rfloor\bmod3,
 \quad 0\leq j<6.
 \label{act21:eq:df-bloque-y-banda}
\end{equation}
Ainsi, \(B_c(t)\) est un mot ordonné de six trits et
\begin{equation}
 u_c(t)=\sum_{j=0}^{5}B_{c,j}(t)3^{5-j}.
 \label{act21:eq:df-elevacion-banda}
\end{equation}
Le mot et son élévation entière sont deux représentations mutuellement récupérables du même bloc émis.

\begin{theorem}[Stabilité d'extraction à toute profondeur]
\label{act21:thm:df-bloque-estable}
On vérifie, pour tout \(n\),
\[
 P_c(729,n)=P_c(3,6n).
\]
Si \(t<n\), alors
\begin{equation}
 u_c(t)=
 \left\lfloor\frac{P_c(729,n)}{729^{n-t-1}}\right\rfloor\bmod729.
 \label{act21:eq:df-bloque-prefijo-largo}
\end{equation}
En conséquence, toute extension du préfixe préserve tous les
blocs précédemment émis et toutes les signatures calculées dessus.
\end{theorem}

\begin{proof}
La première égalité résulte de ce que les deux préfixes se séparent à la même
échelle, \(729^n=3^{6n}\), et du
théorème~\ref{act21:thm:df-emision-correcta}. Pour la seconde, nous appliquons
\eqref{act21:eq:df-truncamiento} aux profondeurs \(n\) et \(t+1\), puis
nous prenons le résidu modulo \(729\). L'extraction ternaire de
\eqref{act21:eq:df-bloque-y-banda} est déterministe. Toute fonction explicite
de ces dix-huit coordonnées conserve, par conséquent, le même résultat en
étendant le préfixe.
\end{proof}

La construction est quantifiée sur \(t\in\mathbb N\). Une table qui
matérialise cent blocs constitue une évaluation finie de cette
construction ; la formule \eqref{act21:eq:df-bloque-prefijo-largo} détermine
exactement sa compatibilité avec toute extension ultérieure.

\begin{example}[Premier bloc conjoint]
L'évaluation des lecteurs régionaux sélectionnés donne les mots
\[
 B_{\mathrm P}(0)=010211,\qquad
 B_{\mathrm E}(0)=201101,\qquad
 B_{\mathrm F}(0)=121200.
\]
Leur élévation par \eqref{act21:eq:df-elevacion-banda} produit
\[
 u_{\mathrm P}(0)=103,\qquad
 u_{\mathrm E}(0)=523,\qquad
 u_{\mathrm F}(0)=450.
\]
Par exemple,
\(201101_3=2\cdot243+27+9+1=523\).
Les six positions, y compris les zéros initiaux, sont conservées dans la
représentation en bande. Le zéro initial de \(010211\) possède une position
définie, bien que l'écriture décimale de \(103\) en soit dépourvue.
\end{example}

\subsection{Registres de transition et reconstruction des élévations linéaires}
\label{act21:subsec:df-registros-transicion}

% Fuentes: deltas/integracion/RegionalTransitionRecords.lean y
% GeneratedTransitionRecords.lean; lean/biblioteca/TPKLifts.lean.
% Los enlaces reconstruidos son R(0)->R(1) y R(2)->R(3).

Sur \(\mathbb F_3\), nous réunissons deux émissions consécutives de chaque région dans la matrice
\begin{equation}
 R(t)=
 \begin{pmatrix}
 B_{\mathrm P}(t)\\ B_{\mathrm P}(t+1)\\
 B_{\mathrm E}(t)\\ B_{\mathrm E}(t+1)\\
 B_{\mathrm F}(t)\\ B_{\mathrm F}(t+1)
 \end{pmatrix}.
 \label{act21:eq:df-registro-matricial}
\end{equation}
Cette définition fixe à la fois l'ordre régional et l'ordre temporel. Les
quatre premiers registres sont
\[
 X_0=R(0),\quad Y_0=R(1),\quad X_1=R(2),\quad Y_1=R(3).
\]
Ses valeurs, obtenues par extraction des préfixes régionaux, sont
\begin{align*}
 X_0&=\begin{pmatrix}
0&1&0&2&1&1\\0&1&2&2&2&2\\2&0&1&1&0&1\\
1&2&1&2&2&1\\1&2&1&2&0&0\\1&1&2&2&0&2
\end{pmatrix},&
 Y_0&=\begin{pmatrix}
0&1&2&2&2&2\\0&1&0&2&1&1\\1&2&1&2&2&1\\
1&0&2&0&1&1\\1&1&2&2&0&2\\1&2&1&0&2&0
\end{pmatrix},\\[1ex]
 X_1&=\begin{pmatrix}
0&1&0&2&1&1\\0&0&2&1&1&1\\1&0&2&0&1&1\\
0&1&2&2&2&2\\1&2&1&0&2&0\\0&1&0&2&1&0
\end{pmatrix},&
 Y_1&=\begin{pmatrix}
0&0&2&1&1&1\\1&1&0&2&2&1\\0&1&2&2&2&2\\
1&0&2&0&1&1\\0&1&0&2&1&0\\0&1&0&2&0&0
\end{pmatrix}.
\end{align*}

\begin{theorem}[Reconstruction unique des deux liens initiaux]
\label{act21:thm:df-elevaciones-reconstruidas}
Les équations \(X_0L_0=Y_0\) et \(X_1L_1=Y_1\) ont, chacune, une
unique solution dans \(M_6(\mathbb F_3)\). Ces solutions sont
\begin{align*}
 L_0&=\begin{pmatrix}
2&2&2&1&2&1\\2&1&2&2&1&1\\1&0&0&1&0&2\\
0&0&1&1&1&0\\2&2&2&0&2&1\\2&1&2&1&0&0
\end{pmatrix},&
 L_1&=\begin{pmatrix}
2&1&2&0&2&2\\1&2&1&1&2&0\\1&2&1&1&1&2\\
0&1&0&0&2&1\\2&0&2&1&0&1\\0&2&2&2&1&1
\end{pmatrix}.
\end{align*}
Les deux matrices sont inversibles.
\end{theorem}

\begin{proof}
Les matrices
\begin{align*}
 A_0&=\begin{pmatrix}
1&1&2&0&1&2\\0&0&1&0&0&1\\2&1&2&1&0&1\\
0&2&0&1&0&2\\2&1&1&0&2&2\\2&1&1&1&1&2
\end{pmatrix},&
 A_1&=\begin{pmatrix}
1&0&1&2&0&0\\0&0&1&1&2&2\\0&1&2&0&1&1\\
1&1&0&2&0&0\\1&1&2&1&1&2\\1&0&0&0&0&2
\end{pmatrix}
\end{align*}
satisfont \(A_iX_i=X_iA_i=I\), par multiplication dans \(\mathbb F_3\).
Par conséquent, les solutions sont imposées par \(L_i=A_iY_i\) ; la
multiplication produit les deux matrices montrées. Leurs inverses sont
\begin{align*}
 L_0^{-1}&=\begin{pmatrix}
1&2&0&2&0&2\\2&0&2&2&0&0\\2&1&2&0&2&0\\
1&0&0&0&2&0\\0&2&1&1&2&0\\2&2&2&2&2&2
\end{pmatrix},&
 L_1^{-1}&=\begin{pmatrix}
2&0&2&1&2&1\\2&1&0&0&2&0\\2&0&2&2&0&2\\
2&0&0&1&1&0\\1&1&1&1&1&0\\2&0&1&2&2&0
\end{pmatrix}.
\end{align*}
La vérification bilatérale \(L_iL_i^{-1}=L_i^{-1}L_i=I\) démontre la
récupération exacte de chaque mot transformé. Tous les produits
utilisent des résidus \(0,1,2\), avec leur interprétation en \(\mathbb F_3\).
\end{proof}

La lecture causale de ce résultat est précise : les émissions régionales produisent \(R(t)\) ; les registres \(R(0),\ldots,R(3)\) déterminent les deux élévations précédentes ; leurs inverses récupèrent les mots transformés.  
La continuité pour tout \(t\) appartient aux lecteurs coinductifs et au théorème ~\ref{act21:thm:df-bloque-estable}. La reconstruction matricielle démontrée ici a pour domaine les deux liens spécifiés dans son énoncé.

\subsection{Raffinement simultané dans les bases ternaire et décimale}
\label{act21:subsec:df-cilindros-mixtos}

% Fuente: deltas/integracion/RegionalCylinderTransition.lean.

Une émission ternaire de six trits et une émission décimale de trois chiffres
constituent deux lectures de bases différentes. Pour conserver leur relation, on
enregistre l'état
\[
 s=(n,k,N,D),
\]
où \(N\) est un préfixe de profondeur \(n\) en base \(729\), et \(D\) est
un préfixe de profondeur \(k\) en base \(1000\). Leurs cylindres
archimédiens sont
\begin{equation}
 I_{729}(s)=\left[\frac{N}{729^n},\frac{N+1}{729^n}\right),\qquad
 I_{1000}(s)=\left[\frac{D}{1000^k},\frac{D+1}{1000^k}\right).
 \label{act21:eq:df-cilindros}
\end{equation}
L'extrémité supérieure semi-ouverte attribue chaque valeur à une unique cellule d'une base et d'une profondeur données.

\begin{definition}[Défauts entiers de compatibilité]
On définit
\begin{align}
 A(s)&=N1000^k-D729^n,\label{act21:eq:df-defecto-inferior}\\
 B(s)&=(N+1)1000^k-D729^n=A(s)+1000^k.
 \label{act21:eq:df-defecto-superior}
\end{align}
L'état est compatible lorsque
\begin{equation}
 A(s)<729^n,\qquad B(s)>0.
 \label{act21:eq:df-compatibilidad}
\end{equation}
\end{definition}

\begin{lemma}[Critère entier d'intersection]
Les cylindres de \eqref{act21:eq:df-cilindros} ont une intersection non vide si et
seulement si \eqref{act21:eq:df-compatibilidad} est vérifiée.
\end{lemma}

\begin{proof}
Deux intervalles semi-ouverts \([a,b)\) et \([c,d)\), avec \(a<b\) et
\(c<d\), se coupent exactement lorsque \(a<d\) et \(c<b\). En appliquant cette
condition à \eqref{act21:eq:df-cilindros} et en multipliant par les dénominateurs
positifs, nous obtenons
\(N1000^k<(D+1)729^n\) et
\(D729^n<(N+1)1000^k\). Ce sont les deux inégalités de
\eqref{act21:eq:df-compatibilidad}.
\end{proof}

\begin{definition}[Mise à jour des profondeurs et des préfixes]
Pour \(0\leq u<729\) et \(0\leq d<1000\), nous définissons
\begin{align*}
 T_u(n,k,N,D)&=(n+1,k,729N+u,D),\\
 T_{u,d}(n,k,N,D)&=(n+1,k+1,729N+u,1000D+d).
\end{align*}
La première opération raffine uniquement le cylindre ternaire. La seconde  
raffine les deux lectures et conserve les deux profondeurs dans l'état.
\end{definition}

\begin{proposition}[Récurrences exactes de frontière]
\label{act21:prop:df-rec-frontera}
Les mises à jour précédentes satisfont
\begin{align}
 A(T_u s)&=729A(s)+u1000^k,\label{act21:eq:df-rec-ternaria}\\
 A(T_{u,d}s)&=729000A(s)+u1000^{k+1}-d729^{n+1}.
 \label{act21:eq:df-rec-conjunta}
\end{align}
\end{proposition}

\begin{proof}
Pour la première identité,
\[
 (729N+u)1000^k-D729^{n+1}
 =729(N1000^k-D729^n)+u1000^k.
\]
Pour la seconde,
\begin{align*}
 &(729N+u)1000^{k+1}-(1000D+d)729^{n+1}\\
 &\qquad=729000(N1000^k-D729^n)
      +u1000^{k+1}-d729^{n+1}.
\end{align*}
Les deux sont des identités entières antérieures à toute approximation réelle.
\end{proof}

\begin{theorem}[Raffinement des cylindres régionaux générés]
\label{act21:thm:df-refinamiento-generado}
Soit
\[
 s_c(n,k)=(n,k,P_c(729,n),P_c(1000,k)).
\]
Alors \(s_c(n,k)\) est compatible quels que soient \(n,k\). De plus,
\begin{align*}
 T_{u_c(n)}s_c(n,k)&=s_c(n+1,k),\\
 T_{u_c(n),d_c(1000,k)}s_c(n,k)&=s_c(n+1,k+1).
\end{align*}
La division euclidienne des préfixes raffinés récupère exactement les
préfixes précédents. Pour tout \(a,b\geq0\),
\[
 \frac{P_c(729,n+a)-P_c(729,n+a)\bmod729^a}{729^a}
 =P_c(729,n),
\]
et l'identité analogue est vérifiée en base \(1000\).
\end{theorem}

\begin{proof}
Le théorème~\ref{act21:thm:df-emision-correcta} place \(v_c\) simultanément dans
les deux intervalles \eqref{act21:eq:df-cilindros} ; le critère d'intersection donne
la compatibilité. Les identités de mise à jour sont
\eqref{act21:eq:df-paso-posicional} dans chaque base. La récupération s'obtient à partir de
\eqref{act21:eq:df-truncamiento}. Par conséquent, la récurrence du défaut et la
récupération du préfixe sont des conséquences d'une même mise à jour.
\end{proof}

\begin{example}[Premier couple de cylindres régionaux]
Pour \(n=k=1\), l'évaluation des trois régions donne
\[
\begin{array}{c|rr|rr}
 c&N&D&A&B\\\hline
 \mathrm P&2290&3141&211&1211\\
 \mathrm E&1981&2718&-422&578\\
 \mathrm F&1179&1618&-522&478
\end{array}
\]
Dans les trois lignes \(A<729\) et \(B>0\). La première colonne de préfixes
se décompose comme \(2290=3\cdot729+103\),
\(1981=2\cdot729+523\) et \(1179=729+450\) : réapparaissent exactement les
blocs de six trits déjà émis. Les signes différents de \(A\) expriment
la position relative des deux extrémités inférieures, sans altérer la
appartenance commune de la valeur régionale.
\end{example}

\subsection{Sélection du bloc suivant par le lecteur régional}

La compatibilité de deux intervalles exprime une relation entre lectures. La détermination du bloc suivant utilise en outre l'intervalle rationnel produit par la région. Cette opération peut se formuler sans éliminer l'histoire de profondeur.

\begin{definition}[Admissibilité de la prolongation régionale]
\(c,n\) étant fixés, soient \(s=729^{n+1}\) et \(j=j_c(s)\). Un résidu
\(u\in\{0,\ldots,728\}\) est admissible lorsque
\begin{equation}
 \lfloor s\ell_c(j)\rfloor
 \leq729P_c(729,n)+u
 \leq\lceil su_c(j)\rceil-1.
 \label{act21:eq:df-admisibilidad}
\end{equation}
\end{definition}

\begin{theorem}[Unicité de la prolongation compatible avec le lecteur]
\label{act21:thm:df-prolongacion-unica}
La condition \eqref{act21:eq:df-admisibilidad} est équivalente à \(u=u_c(n)\).
Par conséquent, il existe exactement une prolongation admissible de chaque région
à chaque profondeur.
\end{theorem}

\begin{proof}
À l'indice de séparation, les deux extrémités entières de
\eqref{act21:eq:df-admisibilidad} coïncident avec \(P_c(729,n+1)\).
La condition est ainsi équivalente à
\(729P_c(729,n)+u=P_c(729,n+1)\).
Par \eqref{act21:eq:df-paso-posicional}, le côté droit est
\(729P_c(729,n)+u_c(n)\). L'annulation du terme commun donne l'égalité des résidus. Réciproquement, cette égalité satisfait la condition.
\end{proof}

La dépendance effective de cette sélection reste complètement déterminée :
région sélectionnée, lecteur rationnel, profondeur, indice de séparation,
préfixe antérieur et résidu suivant. Le défaut \(A\) résume une relation
de frontière ; l'état \((n,k,N,D)\) et les intervalles régionaux conservent
les données employées par la mise à jour.

\subsection{Signature ternaire, retenue, incidence de Witt et horloge de lecture}
\label{act21:subsec:df-firma-conjunta}

% Fuentes: lean/regional/N69RegionalSignature.lean,
% deltas/integracion/JointRegionalFrontier.lean y JointCylinderSignature.lean.

Une profondeur \(t\) étant fixée, écrivons \(B_{c,j}=B_{c,j}(t)\). Pour chaque colonne \(j\), on calcule les entiers et les résidus suivants :
\begin{align}
 S_j&=B_{\mathrm P,j}+B_{\mathrm E,j}+B_{\mathrm F,j},
 &q_j&=S_j\bmod3,
 &c_j&=\left\lfloor S_j/3\right\rfloor,
 \label{act21:eq:df-firma-qc}\\
 a_j&=(B_{\mathrm P,j}+B_{\mathrm E,j}-B_{\mathrm F,j})\bmod3,
 &w_j&=\sum_c\mathbf1_{B_{c,j}\ne0},
 &\bar w_j&=w_j\bmod3.
 \label{act21:eq:df-firma-axial}
\end{align}
Le résidu de l'expression axiale est pris dans les entiers. Une implémentation
avec des naturels le calcule comme
\((B_{\mathrm P,j}+B_{\mathrm E,j}+3-B_{\mathrm F,j})\bmod3\), qui le conserve
exactement.

La matrice d'incidence ternaire utilisée par le lecteur est
\begin{equation}
 W=\begin{pmatrix}
0&1&1&1&1&1\\
1&0&1&1&2&2\\
1&1&0&2&1&2\\
1&1&2&0&2&1\\
1&2&1&2&0&1\\
1&2&2&1&1&0
\end{pmatrix}.
\label{act21:eq:df-matriz-witt}
\end{equation}
Pour chaque région, deux charges résiduelles sont conservées :
\begin{equation}
 v_c^{\mathrm{res}}=\sum_j B_{c,j}\bmod3,
 \qquad
 v_c^{\mathrm W}=\sum_j(B_cW)_j\bmod3.
 \label{act21:eq:df-cargas-residuales}
\end{equation}
La signature complète du lecteur est
\(\mathcal S(B)=(q,a,c,w,\bar w,v^{\mathrm{res}},v^{\mathrm W})\).
Sa structure conserve séparément le reste de colonne, la retenue, le
contraste axial, le poids de support et les deux lectures de charge.

\begin{theorem}[Récupération exacte et bornes de la signature]
\label{act21:thm:df-recuperacion-firma}
Pour chaque colonne et chaque région, on a
\begin{align}
 S_j&=q_j+3c_j,&0\leq q_j,a_j&<3,&0\leq c_j&\leq2,
 \label{act21:eq:df-recuperacion-total}\\
 B_{\mathrm F,j}&=2(q_j+3-a_j)\bmod3,&0\leq w_j&\leq3,
 &0\leq\bar w_j&<3.
 \label{act21:eq:df-recuperacion-eje}
\end{align}
De plus,
\begin{equation}
 v_c^{\mathrm W}
 =\left(\sum_j\bigl((B_cW)_j\bmod3\bigr)\right)\bmod3.
 \label{act21:eq:df-witt-reduccion}
\end{equation}
Les récupérations énoncées déterminent le total de chaque colonne et sa
coordonnée d'auto-échelle ; le registre régional conserve en outre
les trois bandes ordonnées.
\end{theorem}

\begin{proof}
L'identité \eqref{act21:eq:df-recuperacion-total} est la division euclidienne de
\(S_j\) par \(3\). Chaque terme de \(S_j\) est entre \(0\) et \(2\),
de telle sorte que \(0\leq S_j\leq6\) et \(c_j\leq2\). En \(\mathbb F_3\),
la soustraction des deux congruences qui définissent \(q_j\) et \(a_j\) donne
\(q_j-a_j=2B_{\mathrm F,j}\). Comme l'inverse de \(2\) est \(2\), on
obtient \eqref{act21:eq:df-recuperacion-eje} ; le terme \(3\) permet
d'écrire celle-ci avec des nombres naturels. Les bornes de \(w_j\) résultent de sommer trois
indicateurs. Enfin, réduire une somme entière modulo \(3\) équivaut
à réduire chaque terme et à réduire à nouveau la somme, ce qui prouve
\eqref{act21:eq:df-witt-reduccion}.
\end{proof}

\begin{example}[Signature du premier bloc conjoint]
Les trois bandes initiales produisent
\begin{align*}
 q&=(0,0,2,2,1,2),&a&=(1,2,0,1,1,2),\\
 c&=(1,1,0,1,0,0),&w&=(2,2,2,3,1,2),\\
 v^{\mathrm{res}}&=(2,2,0),&v^{\mathrm W}&=(2,1,1).
\end{align*}
La quatrième colonne a un total \(2+1+2=5\) ; sa signature enregistre \(q_3=2\) et \(c_3=1\), par conséquent \(q_3+3c_3=5\).
La récupération axiale donne \(2(2+3-1)\bmod3=2\), précisément la coordonnée d'auto-échelle.
Les images de Witt, réduites composante par composante, sont
\[
 (2,0,2,1,1,2),\qquad
 (0,0,0,2,0,2),\qquad
 (2,1,1,2,1,0).
\]
Sa somme résiduelle par ligne produit les trois charges \((2,1,1)\).
\end{example}

% Recuperación aditiva del propietario N57, con enlace a las pruebas
% de cilindros y firma que ya están incluidas en el manuscrito.
\subsection{Cocycle de retenue et relèvement entier de la signature régionale}
\label{carry27:seccion}

La signature régionale conserve simultanément le résidu et le quotient.
La construction suivante rend explicite la loi algébrique de cette
dernière donnée. APP fournit les évaluations entières sur ses deux
feuillets ; TRIT conserve l'orientation et le régime ; le transport TPK
produit les bandes régionales et met à jour sa mémoire. Dans une
colonne de ces bandes, la division euclidienne détermine la retenue.
La construction utilise les représentants \(\{0,1,2\}\) de
\(\mathbb F_3\) ; la représentation orientée du TRIT conserve
séparément son signe et son quotient de transport.

Soient \((p,e,f)\in\{0,1,2\}^3\) les entrées d'une colonne régionale. Dans cette section, la lettre \(e\) désigne l'entrée ternaire de la région de propagation. Nous définissons
\begin{equation}
\begin{aligned}
q&=[p+e+f]_3,& c&=(p+e+f-q)/3,\\
a&=[p+e-f]_3,& h&=(p+e-f-a)/3,
\end{aligned}
\label{carry27:division}
\end{equation}
où \([z]_3\in\{0,1,2\}\) est le reste de tout entier \(z\).
La variable \(a\) est la coordonnée axiale de cette colonne. La
constante \(\alpha_{\rm HMT}\) conserve sa construction
dodécaphasique postérieure.

\begin{proposition}[Composition locale des retenues]
\label{carry27:cociclo}
Pour \(x,y,z\in\{0,1,2\}\), soient
\[
\kappa_3(x,y)=\frac{x+y-[x+y]_3}{3},\qquad
\beta_3(x,y)=\frac{x-y-[x-y]_3}{3}.
\]
Avec \(r=[p+e]_3\), les deux lois sont
\[
c=\kappa_3(p,e)+\kappa_3(r,f),\qquad
h=\kappa_3(p,e)+\beta_3(r,f).
\]
De plus,
\begin{equation}
\kappa_3(x,y)+\kappa_3([x+y]_3,z)
=\kappa_3(y,z)+\kappa_3(x,[y+z]_3).
\label{carry27:identidad}
\end{equation}
\end{proposition}

\begin{proof}
L'égalité \(p+e=r+3\kappa_3(p,e)\), ajoutée à
\(r+f=q+3\kappa_3(r,f)\), donne la première formule.
En soustrayant \(f\) et en utilisant
\(r-f=a+3\beta_3(r,f)\) on obtient la seconde.
Les deux membres de \eqref{carry27:identidad} sont
\((x+y+z-[x+y+z]_3)/3\). Par conséquent, la composition
préserve le même quotient quelle que soit l'association
de la somme.
\end{proof}

Si \(s_3:\mathbb F_3\to\mathbb Z\) est la section des
représentants, alors
\[
s_3(x)+s_3(y)-s_3(x+y)=3\kappa_3(x,y).
\]
Ceci est le cocycle de l'extension
\(0\to\mathbb Z\xrightarrow{\times3}\mathbb Z\to\mathbb F_3\to0\).
Son apparition provient du relèvement d'une opération résiduelle
à ses valeurs entières. La coordonnée visible et la mémoire du
multiple de trois restent liées par une identité exacte.

\Needspace{10\baselineskip} % S27_LAYOUT_TABLE
\begin{longtable}{ccc|rrrr}
\caption{Loi complète des \(27\) colonnes régionales.
Les valeurs de \(h\) sont présentées comme des entiers.}
\label{carry27:tabla}\\
\toprule
\(p\)&\(e\)&\(f\)&\(q\)&\(c\)&\(a\)&\(h\)\\
\midrule\endfirsthead
\toprule
\(p\)&\(e\)&\(f\)&\(q\)&\(c\)&\(a\)&\(h\)\\
\midrule\endhead
\bottomrule\endfoot
0&0&0&0&0&0&0\\
0&0&1&1&0&2&-1\\
0&0&2&2&0&1&-1\\
0&1&0&1&0&1&0\\
0&1&1&2&0&0&0\\
0&1&2&0&1&2&-1\\
0&2&0&2&0&2&0\\
0&2&1&0&1&1&0\\
0&2&2&1&1&0&0\\
1&0&0&1&0&1&0\\
1&0&1&2&0&0&0\\
1&0&2&0&1&2&-1\\
1&1&0&2&0&2&0\\
1&1&1&0&1&1&0\\
1&1&2&1&1&0&0\\
1&2&0&0&1&0&1\\
1&2&1&1&1&2&0\\
1&2&2&2&1&1&0\\
2&0&0&2&0&2&0\\
2&0&1&0&1&1&0\\
2&0&2&1&1&0&0\\
2&1&0&0&1&0&1\\
2&1&1&1&1&2&0\\
2&1&2&2&1&1&0\\
2&2&0&1&1&1&1\\
2&2&1&2&1&0&1\\
2&2&2&0&2&2&0\\
\end{longtable}

\begin{theorem}[Composante algébrique supplémentaire de la retenue]
\label{carry27:rango}
Sur \(\mathbb F_3\), les fonctions \(q,a\) sont linéaires en
\((p,e,f)\). Les réductions de \(c,h\) ont les représentations polynomiales suivantes :
\begin{align}
\bar c={}&2ef+2ef^2+2e^2f+2pf+2pf^2\notag\\
 &+2pe+pef+2pe^2+2p^2f+2p^2e,
\label{carry27:polinomio-c}\\
\bar h={}&2f^2+ef+2ef^2+e^2f+pf+2pf^2\notag\\
 &+2pe+2pef+2pe^2+p^2f+2p^2e.
\label{carry27:polinomio-h}
\end{align}
Si \(V\) est la matrice d'évaluation de \(1,p,e,f\) sur les
\(27\) colonnes, on a
\[
\operatorname{rank}_{\mathbb F_3}V=4,\qquad
\operatorname{rank}_{\mathbb F_3}(V\mid\bar c)=5,\qquad
\operatorname{rank}_{\mathbb F_3}(V\mid\bar h)=5.
\]
\end{theorem}

\begin{proof}
La substitution des \(27\) triplets du tableau dans les deux
polynômes reproduit \(\bar c,\bar h\). La représentation de
degré au plus deux en chaque variable est unique : un polynôme de
degré au plus deux en une variable qui s'annule aux trois
éléments de \(\mathbb F_3\) est nul ; l'application successive
de cet argument aux trois variables donne l'indépendance
des \(27\) monômes.

Les évaluations en \(0,(1,0,0),(0,1,0),(0,0,1)\) prouvent
l'indépendance de \(1,p,e,f\). La fonction \(\bar c\) vaut
zéro sur ces quatre entrées et vaut un sur \((1,1,1)\) ;
par conséquent, elle reste en dehors de son espace affine d'évaluation.
Pour \(\bar h\), les trois premières entrées imposent les
coefficients constant, de \(p\) et de \(e\) à zéro, et
\(\bar h(0,0,1)=2\) impose le coefficient de \(f\) à deux.
Cette fonction affine vaudrait un sur \((0,0,2)\), tandis
que \(\bar h(0,0,2)=2\). Chaque colonne supplémentaire augmente,
par conséquent, le rang d'une unité.
\end{proof}

La preuve précise quelle information apporte le quotient. La superposition résiduelle produit \(q,a\) ; l'évaluation entière incorpore en outre \(c,h\). Pour \(c\in\{0,1,2\}\), son résidu détermine la valeur entière, et pour \(h\in\{-1,0,1\}\), il la détermine conjointement avec ce choix déclaré de représentants. Les polynômes décrivent les réductions ; la mémoire intégrale conserve également ces relèvements.

\subsubsection{Transport des quotients régionaux}

Soit \(M\) une matrice entière de transport déjà fixée à une étape
du développement, et soient \(x_t^\chi\) ses états de ligne pour
\(\chi\in\{\pi,e,\varphi\}\). Pour chaque coordonnée, on effectue
la division
\[
y_t^\chi=x_t^\chi M,\qquad
u_t^\chi=[y_t^\chi]_3,\qquad
x_{t+1}^\chi=(y_t^\chi-u_t^\chi)/3.
\]
Définissons \(X_t^+=x_t^\pi+x_t^e+x_t^\varphi\) et
\(X_t^\sigma=x_t^\pi+x_t^e-x_t^\varphi\).
Les lois précédentes, appliquées colonne par colonne, donnent
\begin{equation}
\begin{aligned}
X_{t+1}^+&=\frac{X_t^+M-q_t}{3}-c_t,\\
X_{t+1}^\sigma&=\frac{X_t^\sigma M-a_t}{3}-h_t.
\end{aligned}
\label{carry27:transporte}
\end{equation}
En effet, la somme des trois résidus est \(q_t+3c_t\) ; la combinaison signée est \(a_t+3h_t\). Substituer les deux identités dans la division des états démontre \eqref{carry27:transporte}. L'énoncé a pour domaine les transports entiers déclarés. Les émissions régionales et leur continuité conservent les lecteurs coinductifs construits précédemment.

Comme évaluation complète sur une bande de six colonnes,
\[
u^\pi=021020,\quad u^e=001220,\quad u^\varphi=100210
\]
produisent \(q=122120\), \(c=000110\), \(a=222000\) et  
\(h=(-1,0,0,0,1,0)\). L'orientation  
\(\sigma=(1,1,-1)\) possède un résidu double  
\(2\sigma=(2,2,1)\) en \(\mathbb F_3^3\). Cette dernière  
égalité enregistre la convention de signe et se distingue des  
bandes de six colonnes.

\subsubsection{Mémoire de frontière et lecture de l'intervalle}

La loi locale de retenue s'accouple au raffinement cylindrique
déjà démontré. Un exemple montre la fonction de l'intervalle
complet :
\[
\left[\frac1{729},\frac2{729}\right)
\cap
\left[\frac2{1000},\frac3{1000}\right)
=
\left[\frac2{1000},\frac2{729}\right)\ne\varnothing.
\]
En revanche,  
\(\lfloor1000/729\rfloor=1\).  
La lecture de l'extrémité inférieure renvoie le bloc \(1\),  
alors que la compatibilité par intervalles admet également la  
cellule \(2\). L'état de frontière conserve les deux bornes  
nécessaires pour décider de cette compatibilité.

Pour une colonne, fixer \(q\) et \(a\) fixe
\[
f=2(q-a),\qquad p+e=q-f\quad\text{en }\mathbb F_3.
\]
Chacune des neuf fibres de \((q,a)\) contient exactement trois triplets : choisir \(p\) détermine \(e\) et laisse \(f\) fixe. Sur six colonnes, avant d'appliquer des restrictions supplémentaires, la liberté est le produit de ces six fibres et a pour cardinal \(3^6=729\). La fixation se rapporte aux alternatives d'une même frontière. En avançant en profondeur, le registre et sa mémoire changent selon le transport. Ainsi, la retenue locale, les bornes du cylindre et l'incidence régionale agissent conjointement sur la même histoire arithmétique.


\subsection{Horloge entière de lecture décimale}

\begin{definition}[Horloge entière de lecture décimale]
L'indice de transition \(t\) et la profondeur décimale de représentation sont
reliés par
\begin{equation}
 \kappa(t)=\max\{k\in\mathbb N:1000^k\leq3^{6(t+1)}\}.
 \label{act21:eq:df-reloj}
\end{equation}
\end{definition}

\begin{proposition}[Caractérisation et monotonie de l'horloge]
\label{act21:prop:df-reloj}
L'entier \(\kappa(t)\) est déterminé de manière univoque par
\begin{equation}
 1000^{\kappa(t)}\leq729^{t+1}<1000^{\kappa(t)+1}.
 \label{act21:eq:df-reloj-cotas}
\end{equation}
La fonction \(\kappa\) est monotone et satisfait
\(\kappa(20)=\kappa(21)=20\).
\end{proposition}

\begin{proof}
Les puissances de \(1000\) sont strictement croissantes et non bornées ;
il existe donc un unique intervalle de puissances consécutives qui
contient \(729^{t+1}\). La monotonie de cette dernière expression démontre
celle de \(\kappa\). Les deux évaluations finales sont les comparaisons
entières
\[
 1000^{20}\leq729^{21}<1000^{21},\qquad
 1000^{20}\leq729^{22}<1000^{21}.
\]
Ces comparaisons s'effectuent avec des puissances entières exactes.
\end{proof}

L'égalité de profondeur décimale en deux transitions distinctes explique la nécessité de conserver les deux indices : l'horloge de lecture peut rester constante tandis que la bande ternaire avance et que sa signature est mise à jour. L'information chronologique se trouve dans le registre ordonné des transitions.

\subsection{Théorème de composition du bloc, du cylindre et de la signature}

\begin{theorem}[Compatibilité conjointe des lectures régionales]
\label{act21:thm:df-composicion-conjunta}
Pour chaque profondeur \(n\), il existe un unique triplet de blocs
\((u_{\mathrm P},u_{\mathrm E},u_{\mathrm F})\) qui satisfait
l'admissibilité régionale \eqref{act21:eq:df-admisibilidad} dans les trois régions.
Ce triplet est \((u_{\mathrm P}(n),u_{\mathrm E}(n),u_{\mathrm F}(n))\).
Simultanément :
\begin{enumerate}
\item chaque bloc met à jour son cylindre par les récurrences
\eqref{act21:eq:df-rec-ternaria} et \eqref{act21:eq:df-rec-conjunta} ;
\item ses six trits produisent la signature
\(\mathcal S(B_{\mathrm P}(n),B_{\mathrm E}(n),B_{\mathrm F}(n))\);
\item les totaux de colonne se retrouvent comme \(q+3c\), et les charges de Witt sont obtenues à partir des mêmes bandes par
\eqref{act21:eq:df-witt-reduccion} ;
\item tout préfixe de plus grande profondeur renvoie exactement ces
blocs, cylindres tronqués et signatures.
\end{enumerate}
\end{theorem}

\begin{proof}
Nous appliquons le théorème ~\ref{act21:thm:df-prolongacion-unica} à chacun des trois
indices régionaux. L'unicité coordonnée par coordonnée donne l'unicité
conjointe. Nous substituons ces mêmes résidus dans les opérations
\(T_u,T_{u,d}\), ce qui permet d'appliquer le
théorème ~\ref{act21:thm:df-refinamiento-generado}. L'application déterministe de
\eqref{act21:eq:df-bloque-y-banda} produit les bandes utilisées par
\eqref{act21:eq:df-firma-qc}--\eqref{act21:eq:df-cargas-residuales}. Les identités
de récupération sont celles du théorème ~\ref{act21:thm:df-recuperacion-firma}.
Finalement, le théorème ~\ref{act21:thm:df-bloque-estable} et le troncage de
chaque base démontrent la stabilité à plus grande profondeur. À chaque étape, on a conservé l'identité du bloc émis ; aucune des deux
lectures ne requiert de sélectionner un autre bloc.
\end{proof}

La conclusion est une compatibilité opérationnelle entre la contraction
cylindrique et la conservation de l'incidence. Les deux utilisent le même
registre généré, maintiennent les profondeurs et admettent la récupération
entière. Ce résultat est la forme explicite du lien
\[
 \text{bloc régional émis}
 \longrightarrow
 \begin{cases}
 \text{cylindre et frontière entière},\\
 \text{bande, signature, retenue et charge de Witt}.
 \end{cases}
\]
Le développement ultérieur de l'incidence et des champs emploie la deuxième lecture, conservant l'origine commune avec la première.
